\chapter{分析化学}
\section{陽イオンの系統分析}
\label{theory:analytical:separation}
\subsection{背景と基本概念}
混合溶液の成分を、一度に名前だけで判定するのでなく、条件を変えながら分離する。
沈殿ができた段階でその成分はろ液から除かれる。後の手順を解くときは、既に除いたイオンを候補に戻さない。
\subsection{物理専攻者への接続}
沈殿の選択性は平衡条件と溶解度で説明できる。イオン積と溶解度積を比較し、pHが反応種の濃度へ与える影響を考える。
硫化水素の酸解離を二段で扱う近似では、条件を満たす範囲で$[\mathrm{S}^{2-}]\propto[\mathrm{H}^+]^{-2}$となり、酸性と塩基性で硫化物の沈殿対象が変わる。
\subsection{過去問の条件}
H26の例は塩化物、酸性下の硫化物、水酸化物、アンモニア性条件の硫化物を順に分離する。
原典の手順IVを手順Iと読むと対象が変わる。実験条件の詳細は教材の簡略化から推測で補わない\cite{qualitative_stanford}。

\section{観測することと分けること}\label{theory:analysis:separation}
分析では、試料から何を知るかと、どの操作で成分を分けるかを区別する。NMRは磁場中の核スピンの共鳴を測り、通常は試料を成分別に回収しない。遠心分離は沈降、電気泳動は電場中の移動度、クロマトグラフィーは固定相と移動相の間の保持の差を利用する。
MSは作ったイオンを質量電荷比で選別するため装置内部の分離を含むが、元の中性混合物を大量に分取する方法という意味とは異なる\cite{analysis49_ms_ref}。原典の「成分分離」の広さを確認する。

\section{分子量はどの平均か}\label{theory:analysis:molar}
分子量の違う鎖が共存する高分子では、一つの値を述べるだけでは分布が分からない。分子数 \(N_i\)、分子量 \(M_i\) に対し
\[
 M_n=\frac{\sum N_iM_i}{\sum N_i},\qquad
 M_w=\frac{\sum N_iM_i^2}{\sum N_iM_i}.
\]
粒子数に依存する浸透圧は数平均に対応する。溶質質量濃度 \(c\)、絶対温度 \(T\)、気体定数 \(R\) を使い、非電解質高分子の希薄溶液では
\[
 \frac{\Pi}{c}=\frac{RT}{M_n}+O(c).
\]
濃度を変えて零へ外挿し、相互作用項を分けて切片を使う。
GPCは溶液中の大きさに応じた細孔への侵入の差を使う。標準物質による保持の校正から相対的な分子量と分布も測定できる\cite{analysis49_gpc_ref}。
鎖形状や溶媒の違いは校正へ影響する。「絶対測定」を明記しない問いで、GPCを分子量測定でないと排除しない。

\section{略称を励起・分離・検出へ分解する}\label{theory:analysis:acronyms}
ESRはElectron Spin Resonance、EPRはElectron Paramagnetic Resonanceで不対電子を対象とする。核磁気共鳴NMRとは対象が違う\cite{analysis49_esr_ref}。
HPLCの通常名はHigh Performance Liquid ChromatographyでPrecisionではない。
MALDIはMatrix-Assisted Laser Desorption/Ionizationでイオン化、TOFは飛行時間型の質量選別を表す\cite{analysis49_hplc_ref,analysis49_maldi_ref}。
CARSはCoherent Anti-Stokes Raman Spectroscopy、ESCAはElectron Spectroscopy for Chemical Analysisである\cite{analysis49_cars_ref,analysis49_xps_ref}。
FTIRのFTはFourier Transformで、核分裂飛跡と区別する。EPMAは電子線による局所分析、NAAは中性子放射化分析、GC/MSは気体クロマトグラフィーと質量分析の組合せである\cite{analysis49_epma_ref,analysis49_naa_ref}。
装置名に含まれる「何を照射するか」「何で分けるか」「何を観測するか」を別々に読む。

\section{吸光・発光と濃度}\label{theory:analysis:beer}
光強度の透過率 \(T_{\rm opt}=I/I_0\)、吸光度 \(A=-\log_{10}T_{\rm opt}\) を定める。
同じ化学種、波長、光路長 \(l\) を保つ希薄溶液で
\[
 A=\varepsilon cl,\qquad \log_{10}T_{\rm opt}=-\varepsilon cl.
\]
透過率自体は濃度に線形ではないが、その対数は負の係数で比例する\cite{analysis49_beer_ref}。濃度で会合・解離が変わる試料や強い散乱、迷光を無条件に同じ式へ入れない。
炎光分析は励起した原子が低いエネルギー状態へ戻る発光を使う。UV--Visは電子遷移、IRは主に振動遷移という励起対象の違いを押さえる。

\section{温度計と天秤にも測定条件がある}\label{theory:analysis:metrology}
温度計の感温部を入れれば常に真の温度を得るわけではない。全浸没型液体温度計と部分浸没型は校正条件が違い、後者では露出柱の温度も関係する\cite{analysis49_thermo_ref}。「水銀柱全部が同温度」は全型を網羅する規則ではない。
天秤では空気の浮力により概略の有効力が \(mg-\rho_{\rm air}Vg\) になる。同質量でも密度の違う分銅と試料は体積が違う。空気密度を温度・圧力・湿度から求め、必要なら補正する\cite{analysis49_buoy_ref}。
物理専攻者は力の釣合いへ戻って誤差の向きを判断できる。高精度測定での補正と、通常の秤量で無視する近似の大きさを区別する。

\section{内殻を測っても結合状態を調べられる}\label{theory:analysis:xray}
内殻電子を励起・電離すると空孔ができる。別の電子がその空孔へ移ると元素に特徴的な蛍光X線を出し得る。入射光子は吸収端より高エネルギー、すなわち短波長であることが必要になる。XRFの複数ピークから複数元素を分析できる\cite{analysis49_xrf_ref}。
XPSでは放出した光電子の運動エネルギーから結合エネルギーを求める。
\[
 E_{\rm bind}=h\nu-E_{\rm kin}-\phi
\]
ここで \(\phi\) は装置の仕事関数等の基準補正である。化学状態で内殻結合エネルギーに差が生じるので「内殻だから化学結合に鈍感」としない\cite{analysis49_xps_ref}。
X線発光にも高分解能で電子状態の差を扱う方法がある\cite{analysis49_sxes_ref}。微小シフトがあることと、任意の通常装置で必ずそれを分解できることは違う。

\section{分離と検出と校正}\label{theory:analysis:chrom}
保持値は固定相・移動相・温度等に依存する。同条件で標準と一致すれば候補を推定できるが、共溶出があり保持一致だけでは唯一の同定とはならない。MS検出器は質量スペクトルという追加情報を得られる\cite{analysis49_library_ref}。
外部標準の検量線では既知量と信号の関係を作る。内標準法は対象と分離される物質を加え、応答比で注入量等の変動を補正する。標準添加は試料の分取へ既知量を加え、そのマトリックスで応答を測る\cite{analysis49_chrom_ref}。
全検出器が一律に線形応答というわけではない。例えばFPD硫黄モードの非線形性がある。検出器の機構と検量線の適合を確認する。

\section{分離操作と共沸の例外}\label{theory:analysis:operations}
蒸留は気液の組成差を使うが、共沸では \(y=x\) となる。単に沸点の違う二成分だから必ず気液組成が違うとは言えない\cite{nptel_azeotrope}。
吸収は気体成分を液体内部へ溶かす操作、吸着は表面や細孔へ保持する操作、液液抽出は非混和二相間の分配を使う操作である。
逆浸透は膜を挟む浸透圧差に逆らう圧力差を与える。水を優先的に透過させる代表的な脱塩法だが、実膜の全溶質阻止率を常に100\%としない。理想膜での原理と実際の透過性能を分ける。

\section{IRは伸縮だけではない}\label{theory:analysis:ir}
二原子の近似では \(\omega\simeq\sqrt{k/\mu_{\rm red}}\) で、力定数と換算質量が振動数を決める。多原子分子には伸縮、結合角を変える変角、面外変形等の多数の正規振動がある。振動に伴って双極子が変わるものがIR吸収へ現れる\cite{analysis49_ir_ref,atomic_ir_raman}。
官能基が同じなら似た吸収領域を持つことが多いが、共役・水素結合・分子環境で変わり得る。Ramanの分極率変化という条件と、IRの双極子変化を交換しない。

\section{発光・解離・軌道エネルギー}\label{theory:analysis:photophysics}
通常の蛍光は同じスピン多重度の励起状態からの発光、りん光は典型的に三重項から一重項への遷移である。系間交差、無放射緩和、発光を一つの過程へまとめない。核は電子に比べて重く、電子遷移では核配置をほぼ固定した垂直遷移と振動波動関数の重なりを扱う。
前期解離は束縛的な状態と解離性の別状態の結合により生じ得る非放射解離である。自身の解離極限を超えた直接解離の単なる別名ではない\cite{analysis49_pred_ref}。
Hartree--Fockで他の軌道を凍結して一電子を取り去るKoopmans近似では
\[
 I_i\simeq-\varepsilon_i.
\]
軌道緩和と電子相関の変化を省く近似であり、全固体・全分子の観測値への厳密な等式ではない\cite{analysis49_koop_ref}。
発光の種類には遅延蛍光等もあり、過去問の典型的な図式と全過程の一般論を区別する。

\section{学習順序と原典の疑義}
まず対象を原子核・電子・振動・元素・分子量へ分け、次に分離の原理と検出信号を整理する。最後に検量線・測定環境・成立条件を確認する。公式の番号を保存しつつ、GPC、検出器の全称、温度計の校正などの限定不足は問題別の未解決記録へ残す。

\section{原典と全五肢の解説初稿}
\section{ESRを核磁気共鳴と取り違えた略称を判断する。}

\subsection{平成24年度 問27}
\label{q:h24-reg-applied-27}

\noindent\textbf{原典：} sources/H24/17\_applied\_science.pdf、PDF物理8頁。

\textbf{公式正答：} 4（single）。

\paragraph{問題の要旨} ESRを核磁気共鳴と取り違えた略称を判断する。

\Needspace{4\baselineskip}\paragraph{確認した条件・記述（独自要約）}\mbox{}\par

\noindent\textbf{原典確認：} S7\_049の26原典画像で全五肢・否定・添字を確認。\par

\noindent\textbf{対象：} ESRを核磁気共鳴と取り違えた略称を判断する。\par

\paragraph{問われる核心} ESRを核磁気共鳴と取り違えた略称を判断する。

\paragraph{必要知識} 磁気共鳴で対象となるスピンを区別する。ESR/EPRは不対電子、NMRはスピンを持つ原子核の共鳴を扱う。電子と核を交換した4が不適切である。

\paragraph{考え方・正答を選ぶ理由（査読待ち初稿）} 磁気共鳴で対象となるスピンを区別する。ESR/EPRは不対電子、NMRはスピンを持つ原子核の共鳴を扱う。電子と核を交換した4が不適切である。

\Needspace{4\baselineskip}\paragraph{選択肢の検討（査読待ち初稿）}\mbox{}\par

\noindent\textbf{1：通常の条件で適切} XRF／蛍光X線。X-ray Fluorescenceである。\par

\noindent\textbf{2：通常の条件で適切} SEM／走査型電子顕微鏡。Scanning Electron Microscopeである。\par

\noindent\textbf{3：通常の条件で適切} ICP-MS／誘導結合プラズマ質量分析。Inductively Coupled Plasma Mass Spectrometryで、プラズマで作ったイオンを質量分析する。\par

\noindent\textbf{4：不適切} ESR／核磁気共鳴。ESRはElectron Spin Resonance。核磁気共鳴はNMRである。\par

\noindent\textbf{5：通常の条件で適切} FTIR／フーリエ変換赤外分光。Fourier Transform Infraredで、干渉信号から赤外スペクトルを求める。\par

\paragraph{再利用する考え方} 測定対象と分離対象、定義の分母、絶対量と相対校正、全称と成立条件を分けて判断する。

\paragraph{共通理論} \hyperref[theory:analysis:acronyms]{第\ref*{theory:analysis:acronyms}節}

\par 参考：\cite{analysis49_esr_ref}。

\par 難易度 3（校準初稿）、優先度未評価。

\subsection{平成30年度 問13}
\label{q:h30-reg-applied-13}

\noindent\textbf{原典：} sources/H30/17\_applied\_science.pdf、PDF物理4頁。

\textbf{公式正答：} 2（single）。

\paragraph{問題の要旨} HPLCのPerformanceとPrecisionを区別する。

\Needspace{4\baselineskip}\paragraph{確認した条件・記述（独自要約）}\mbox{}\par

\noindent\textbf{原典確認：} S7\_049の26原典画像で全五肢・否定・添字を確認。\par

\noindent\textbf{対象：} HPLCのPerformanceとPrecisionを区別する。\par

\paragraph{問われる核心} HPLCのPerformanceとPrecisionを区別する。

\paragraph{必要知識} HPLCのPを高精度のPrecisionへ交換した肢が不適切。通常はHigh Performance Liquid Chromatographyである。他の略称は観測又はイオン化機構を表す。MALDIはイオン源、TOF等は質量分析器であり、それぞれの機能を分けて覚える。

\paragraph{考え方・正答を選ぶ理由（査読待ち初稿）} HPLCのPを高精度のPrecisionへ交換した肢が不適切。通常はHigh Performance Liquid Chromatographyである。他の略称は観測又はイオン化機構を表す。MALDIはイオン源、TOF等は質量分析器であり、それぞれの機能を分けて覚える。

\Needspace{4\baselineskip}\paragraph{選択肢の検討（査読待ち初稿）}\mbox{}\par

\noindent\textbf{1：通常の条件で適切} EPRの正式名。Electron Paramagnetic Resonance。電子常磁性共鳴でESRと同じ対象を扱う。\par

\noindent\textbf{2：不適切} HPLCをHigh Precisionとする。現在の通常名はHigh Performance Liquid Chromatography。高精度のPrecisionではない。\par

\noindent\textbf{3：通常の条件で適切} ESCAの正式名。Electron Spectroscopy for Chemical Analysis。XPSによる元素・化学状態の分析に用いられる。\par

\noindent\textbf{4：通常の条件で適切} CARSの正式名。Coherent Anti-Stokes Raman Spectroscopy。コヒーレントな反Stokes Raman信号を観測する。\par

\noindent\textbf{5：通常の条件で適切} MALDIの正式名。Matrix-Assisted Laser Desorption/Ionization。マトリックス支援レーザー脱離イオン化を表す。\par

\paragraph{再利用する考え方} 測定対象と分離対象、定義の分母、絶対量と相対校正、全称と成立条件を分けて判断する。

\paragraph{共通理論} \hyperref[theory:analysis:acronyms]{第\ref*{theory:analysis:acronyms}節}

\par 参考：\cite{analysis49_hplc_ref,analysis49_esr_ref,analysis49_maldi_ref,analysis49_cars_ref}。

\paragraph{関連問題} \hyperref[q:r06-reg-applied-14]{R06-REG-APPLIED-14}

\par 難易度 3（校準初稿）、優先度未評価。

\subsection{令和6年度 問14}
\label{q:r06-reg-applied-14}

\noindent\textbf{原典：} sources/R06/17\_applied\_science.pdf、PDF物理4頁。

\textbf{公式正答：} 4（single）。

\paragraph{問題の要旨} HPLCのPerformanceとPrecisionを区別する。

\Needspace{4\baselineskip}\paragraph{確認した条件・記述（独自要約）}\mbox{}\par

\noindent\textbf{原典確認：} S7\_049の26原典画像で全五肢・否定・添字を確認。\par

\noindent\textbf{対象：} HPLCのPerformanceとPrecisionを区別する。\par

\paragraph{問われる核心} HPLCのPerformanceとPrecisionを区別する。

\paragraph{必要知識} HPLCのPを高精度のPrecisionへ交換した肢が不適切。通常はHigh Performance Liquid Chromatographyである。他の略称は観測又はイオン化機構を表す。MALDIはイオン源、TOF等は質量分析器であり、それぞれの機能を分けて覚える。

\paragraph{考え方・正答を選ぶ理由（査読待ち初稿）} HPLCのPを高精度のPrecisionへ交換した肢が不適切。通常はHigh Performance Liquid Chromatographyである。他の略称は観測又はイオン化機構を表す。MALDIはイオン源、TOF等は質量分析器であり、それぞれの機能を分けて覚える。

\Needspace{4\baselineskip}\paragraph{選択肢の検討（査読待ち初稿）}\mbox{}\par

\noindent\textbf{1：通常の条件で適切} CARSの正式名。Coherent Anti-Stokes Raman Spectroscopy。コヒーレントな反Stokes Raman信号を観測する。\par

\noindent\textbf{2：通常の条件で適切} EPRの正式名。Electron Paramagnetic Resonance。電子常磁性共鳴でESRと同じ対象を扱う。\par

\noindent\textbf{3：通常の条件で適切} ESCAの正式名。Electron Spectroscopy for Chemical Analysis。XPSによる元素・化学状態の分析に用いられる。\par

\noindent\textbf{4：不適切} HPLCをHigh Precisionとする。現在の通常名はHigh Performance Liquid Chromatography。高精度のPrecisionではない。\par

\noindent\textbf{5：通常の条件で適切} MALDIの正式名。Matrix-Assisted Laser Desorption/Ionization。マトリックス支援レーザー脱離イオン化を表す。\par

\paragraph{再利用する考え方} 測定対象と分離対象、定義の分母、絶対量と相対校正、全称と成立条件を分けて判断する。

\paragraph{共通理論} \hyperref[theory:analysis:acronyms]{第\ref*{theory:analysis:acronyms}節}

\par 参考：\cite{analysis49_hplc_ref,analysis49_esr_ref,analysis49_maldi_ref,analysis49_cars_ref}。

\paragraph{関連問題} \hyperref[q:h30-reg-applied-13]{H30-REG-APPLIED-13}

\par 難易度 3（校準初稿）、優先度未評価。

\subsection{令和6年度 問19}
\label{q:r06-reg-applied-19}

\noindent\textbf{原典：} sources/R06/17\_applied\_science.pdf、PDF物理6頁。

\textbf{公式正答：} 1（single）。

\paragraph{問題の要旨} FTIRとフィッショントラックを区別する。

\Needspace{4\baselineskip}\paragraph{確認した条件・記述（独自要約）}\mbox{}\par

\noindent\textbf{原典確認：} S7\_049の26原典画像で全五肢・否定・添字を確認。\par

\noindent\textbf{対象：} FTIRとフィッショントラックを区別する。\par

\paragraph{問われる核心} FTIRとフィッショントラックを区別する。

\paragraph{必要知識} フーリエ変換赤外というFTIRと、年代測定等で使う核分裂飛跡を区別する。EPMAは電子線、XRFはX線、NAAは中性子を試料へ与える。GC/MSは分離器と検出分析器の組合せである。

\paragraph{考え方・正答を選ぶ理由（査読待ち初稿）} フーリエ変換赤外というFTIRと、年代測定等で使う核分裂飛跡を区別する。EPMAは電子線、XRFはX線、NAAは中性子を試料へ与える。GC/MSは分離器と検出分析器の組合せである。

\Needspace{4\baselineskip}\paragraph{選択肢の検討（査読待ち初稿）}\mbox{}\par

\noindent\textbf{1：不適切} FTIR／フィッショントラック。FTIRはFourier Transform Infrared。フィッショントラックは核分裂飛跡を用いる別方法。\par

\noindent\textbf{2：通常の条件で適切} EPMA／電子線マイクロアナライザー。Electron Probe Microanalyzerで局所に電子線を照射し特性X線等を測る。\par

\noindent\textbf{3：通常の条件で適切} GC/MS／ガスクロマトグラフ質量分析。GCで分離した成分をMSで測定する複合方法。\par

\noindent\textbf{4：通常の条件で適切} NAA／中性子放射化分析。Neutron Activation Analysis。中性子照射後に生じる放射性核種等の信号から元素を分析する。\par

\noindent\textbf{5：通常の条件で適切} XRF／蛍光X線。X-ray Fluorescenceである。\par

\paragraph{再利用する考え方} 測定対象と分離対象、定義の分母、絶対量と相対校正、全称と成立条件を分けて判断する。

\paragraph{共通理論} \hyperref[theory:analysis:acronyms]{第\ref*{theory:analysis:acronyms}節}

\par 参考：\cite{analysis49_ir_ref,analysis49_epma_ref,analysis49_naa_ref}。

\par 難易度 3（校準初稿）、優先度未評価。

\section{クロマトグラフィーの保持・検出・検量線を区別する。}

\subsection{令和5年度 問22}
\label{q:r05-reg-applied-22}

\noindent\textbf{原典：} sources/R05/17\_applied\_science.pdf、PDF物理8頁。

\textbf{公式正答：} 4（single）。

\par\textbf{検討未確定：} 公式4を保持。肢1はMS等の成分識別情報とFPD硫黄の非線形応答に反例がある。一般濃度検出器という限定が原典になく、全五肢の排他性を査読。

\paragraph{問題の要旨} クロマトグラフィーの保持・検出・検量線を区別する。

\Needspace{4\baselineskip}\paragraph{確認した条件・記述（独自要約）}\mbox{}\par

\noindent\textbf{原典確認：} S7\_049の26原典画像で全五肢・否定・添字を確認。\par

\noindent\textbf{対象：} クロマトグラフィーの保持・検出・検量線を区別する。\par

\paragraph{問われる核心} クロマトグラフィーの保持・検出・検量線を区別する。

\paragraph{必要知識} 内標準と対象成分のピークは分離される必要がある。比を取る二信号を分けて測れない4が公式の不適切肢である。ただし1も検出器を全称化するとMSや非線形検出器に合わない。2の保持一致は同一性の推定であって単独の厳密な証明ではない。

\paragraph{考え方・正答を選ぶ理由（査読待ち初稿）} 内標準と対象成分のピークは分離される必要がある。比を取る二信号を分けて測れない4が公式の不適切肢である。ただし1も検出器を全称化するとMSや非線形検出器に合わない。2の保持一致は同一性の推定であって単独の厳密な証明ではない。

\Needspace{4\baselineskip}\paragraph{選択肢の検討（査読待ち初稿）}\mbox{}\par

\noindent\textbf{1：条件付き・解釈留保} 検出器は常に比例応答で成分同定不可。一般の濃度検出器の略説には対応するがMS検出器は質量スペクトル等を得る。FPD硫黄モードのように応答が量の二乗となる方式もあり、検出器全ての性質ではない。\par

\noindent\textbf{2：条件付き・解釈留保} 保持値一致で同一物質と推定。同条件なら同定の手掛かりになるが、異なる成分が共溶出することもある。推定と確証を区別する。\par

\noindent\textbf{3：通常の条件で適切} 絶対検量線は量とピーク面積の関係。既知濃度の標準試料から応答曲線を作り未知試料へ適用する。名称に反して測定値が校正不要という意味ではない。\par

\noindent\textbf{4：不適切} 内標準に重なるピークを選ぶ。各対象ピークから分離される内標準を選び、面積比等で変動を補正する。重なると双方の信号を独立に測れない。\par

\noindent\textbf{5：通常の条件で適切} 標準添加でマトリックス影響を考慮。同じ試料の分取へ既知量を添加することで試料環境での応答を校正する。直線性・加算性等は確認する。\par

\paragraph{再利用する考え方} 測定対象と分離対象、定義の分母、絶対量と相対校正、全称と成立条件を分けて判断する。

\paragraph{共通理論} \hyperref[theory:analysis:chrom]{第\ref*{theory:analysis:chrom}節}

\par 参考：\cite{analysis49_chrom_ref,analysis49_library_ref,analysis49_hplc_ref}。

\par 難易度 4（校準初稿）、優先度未評価。

\section{IRの伸縮と変角を区別し分光法の対象を対応させる。}

\subsection{令和5年度 問15}
\label{q:r05-reg-applied-15}

\noindent\textbf{原典：} sources/R05/17\_applied\_science.pdf、PDF物理5頁。

\textbf{公式正答：} 2（single）。

\par\textbf{検討未確定：} モル吸光係数を構造だけの無条件定数とはしない。波長・溶媒・化学種等を指定する条件付きの説明として保持。

\paragraph{問題の要旨} IRの伸縮と変角を区別し分光法の対象を対応させる。

\Needspace{4\baselineskip}\paragraph{確認した条件・記述（独自要約）}\mbox{}\par

\noindent\textbf{原典確認：} S7\_049の26原典画像で全五肢・否定・添字を確認。\par

\noindent\textbf{対象：} IRの伸縮と変角を区別し分光法の対象を対応させる。\par

\paragraph{問われる核心} IRの伸縮と変角を区別し分光法の対象を対応させる。

\paragraph{必要知識} IRは伸縮だけでなく変角・曲げ等の振動も観測できる。したがってのみという限定の肢を選ぶ。NMRは核スピン、UV–Visは電子遷移を対象とする。官能基の吸収領域は近似的な構造指標であり、水素結合や共役で変わり得る。

\paragraph{考え方・正答を選ぶ理由（査読待ち初稿）} IRは伸縮だけでなく変角・曲げ等の振動も観測できる。したがってのみという限定の肢を選ぶ。NMRは核スピン、UV–Visは電子遷移を対象とする。官能基の吸収領域は近似的な構造指標であり、水素結合や共役で変わり得る。

\Needspace{4\baselineskip}\paragraph{選択肢の検討（査読待ち初稿）}\mbox{}\par

\noindent\textbf{1：通常の条件で適切} NMRは核スピンとラジオ波の共鳴。静磁場中でスピンを持つ核のエネルギー差に対応する高周波を用いる。\par

\noindent\textbf{2：不適切} IRは伸縮運動のみ。結合角を変える変角・曲げ振動等も双極子を変えれば吸収する。\par

\noindent\textbf{3：通常の条件で適切} UV–Visは電子遷移。光子吸収で分子の電子状態を励起する。\par

\noindent\textbf{4：条件付き・解釈留保} 強度と構造固有のモル吸光係数。εは化合物の構造に関係するが、波長・溶媒・温度・化学種状態等も指定する必要がある。\par

\noindent\textbf{5：通常の条件で適切} 同じ官能基は近い吸収領域。局所的な結合の力定数と換算質量が似るため。共役・水素結合・環境によるずれはある。\par

\paragraph{再利用する考え方} 測定対象と分離対象、定義の分母、絶対量と相対校正、全称と成立条件を分けて判断する。

\paragraph{共通理論} \hyperref[theory:analysis:ir]{第\ref*{theory:analysis:ir}節}

\par 参考：\cite{analysis49_ir_ref,analysis49_beer_ref}。

\paragraph{関連問題} \hyperref[q:r07-reg-applied-14]{R07-REG-APPLIED-14}

\par 難易度 4（校準初稿）、優先度未評価。

\subsection{令和7年度 問14}
\label{q:r07-reg-applied-14}

\noindent\textbf{原典：} sources/R07/17\_applied\_science.pdf、PDF物理5頁。

\textbf{公式正答：} 4（single）。

\par\textbf{検討未確定：} モル吸光係数を構造だけの無条件定数とはしない。波長・溶媒・化学種等を指定する条件付きの説明として保持。

\paragraph{問題の要旨} IRの伸縮と変角を区別し分光法の対象を対応させる。

\Needspace{4\baselineskip}\paragraph{確認した条件・記述（独自要約）}\mbox{}\par

\noindent\textbf{原典確認：} S7\_049の26原典画像で全五肢・否定・添字を確認。\par

\noindent\textbf{対象：} IRの伸縮と変角を区別し分光法の対象を対応させる。\par

\paragraph{問われる核心} IRの伸縮と変角を区別し分光法の対象を対応させる。

\paragraph{必要知識} IRは伸縮だけでなく変角・曲げ等の振動も観測できる。したがってのみという限定の肢を選ぶ。NMRは核スピン、UV–Visは電子遷移を対象とする。官能基の吸収領域は近似的な構造指標であり、水素結合や共役で変わり得る。

\paragraph{考え方・正答を選ぶ理由（査読待ち初稿）} IRは伸縮だけでなく変角・曲げ等の振動も観測できる。したがってのみという限定の肢を選ぶ。NMRは核スピン、UV–Visは電子遷移を対象とする。官能基の吸収領域は近似的な構造指標であり、水素結合や共役で変わり得る。

\Needspace{4\baselineskip}\paragraph{選択肢の検討（査読待ち初稿）}\mbox{}\par

\noindent\textbf{1：通常の条件で適切} NMRは核スピンとラジオ波の共鳴。静磁場中でスピンを持つ核のエネルギー差に対応する高周波を用いる。\par

\noindent\textbf{2：通常の条件で適切} UV–Visは電子遷移。光子吸収で分子の電子状態を励起する。\par

\noindent\textbf{3：条件付き・解釈留保} 強度と構造固有のモル吸光係数。εは化合物の構造に関係するが、波長・溶媒・温度・化学種状態等も指定する必要がある。\par

\noindent\textbf{4：不適切} IRは伸縮運動のみ。結合角を変える変角・曲げ振動等も双極子を変えれば吸収する。\par

\noindent\textbf{5：通常の条件で適切} 同じ官能基は近い吸収領域。局所的な結合の力定数と換算質量が似るため。共役・水素結合・環境によるずれはある。\par

\paragraph{再利用する考え方} 測定対象と分離対象、定義の分母、絶対量と相対校正、全称と成立条件を分けて判断する。

\paragraph{共通理論} \hyperref[theory:analysis:ir]{第\ref*{theory:analysis:ir}節}

\par 参考：\cite{analysis49_ir_ref,analysis49_beer_ref}。

\paragraph{関連問題} \hyperref[q:r05-reg-applied-15]{R05-REG-APPLIED-15}

\par 難易度 4（校準初稿）、優先度未評価。

\section{分析信号と水銀温度計・天秤の環境依存を比較する。}

\subsection{平成29年度 問12}
\label{q:h29-reg-applied-12}

\noindent\textbf{原典：} sources/H29/17\_applied\_science.pdf、PDF物理4頁。

\textbf{公式正答：} 5（single）。

\par\textbf{検討未確定：} 水銀温度計の全浸没型という条件が原典に明記されていない。部分浸没型の校正まで全水銀柱同温と一般化しない。Beer則にも化学種・波長・散乱等の成立条件がある。

\paragraph{問題の要旨} 分析信号と水銀温度計・天秤の環境依存を比較する。

\Needspace{4\baselineskip}\paragraph{確認した条件・記述（独自要約）}\mbox{}\par

\noindent\textbf{原典確認：} S7\_049の26原典画像で全五肢・否定・添字を確認。\par

\noindent\textbf{対象：} 分析信号と水銀温度計・天秤の環境依存を比較する。\par

\paragraph{問われる核心} 分析信号と水銀温度計・天秤の環境依存を比較する。

\paragraph{必要知識} 天秤で下向きの有効力は概略mg−ρair Vg。空気密度は温度・圧力・湿度に依存し、試料と分銅の体積差によって浮力補正が必要となる。全く変わらない5を選ぶ。Beer則A=−log10 T=εclなのでH29のlog透過度も負の比例であり、R02の吸光度も正の比例となる。

\paragraph{考え方・正答を選ぶ理由（査読待ち初稿）} 天秤で下向きの有効力は概略mg−ρair Vg。空気密度は温度・圧力・湿度に依存し、試料と分銅の体積差によって浮力補正が必要となる。全く変わらない5を選ぶ。Beer則A=−log10 T=εclなのでH29のlog透過度も負の比例であり、R02の吸光度も正の比例となる。

\Needspace{4\baselineskip}\paragraph{選択肢の検討（査読待ち初稿）}\mbox{}\par

\noindent\textbf{1：通常の条件で適切} 炎中の原子励起と発光。高い状態から低い状態への遷移による特性発光を測る。\par

\noindent\textbf{2：通常の条件で適切} 透過度の対数の濃度比例。固定波長・光路長でBeer則A=−log10T=εcl。logTも係数の符号が負の比例関係となる。\par

\noindent\textbf{3：条件付き・解釈留保} 水銀柱全部が同温度なら正しい。全浸没用温度計ではその条件に対応する。一方、部分浸没型は指定浸没位置と露出柱の校正条件があり、全型に無条件の規則ではない。\par

\noindent\textbf{4：通常の条件で適切} XRF波長が化学結合状態で変わり得る。内殻の信号にも周辺電子状態に応じた小さいエネルギー差があり、高分解能では化学シフトを扱える。\par

\noindent\textbf{5：不適切} 天秤は温度・圧力・湿度で変化しない。空気密度と浮力が変わる。試料と分銅の密度差があると見かけの比較が変わる。\par

\paragraph{再利用する考え方} 測定対象と分離対象、定義の分母、絶対量と相対校正、全称と成立条件を分けて判断する。

\paragraph{共通理論} \hyperref[theory:analysis:metrology]{第\ref*{theory:analysis:metrology}節}

\par 参考：\cite{analysis49_beer_ref,analysis49_thermo_ref,analysis49_buoy_ref,analysis49_sxes_ref}。

\paragraph{関連問題} \hyperref[q:r02-reg-applied-11]{R02-REG-APPLIED-11}

\par 難易度 4（校準初稿）、優先度未評価。

\subsection{令和2年度 問11}
\label{q:r02-reg-applied-11}

\noindent\textbf{原典：} sources/R02/17\_applied\_science.pdf、PDF物理4頁。

\textbf{公式正答：} 5（single）。

\par\textbf{検討未確定：} 水銀温度計の全浸没型という条件が原典に明記されていない。部分浸没型の校正まで全水銀柱同温と一般化しない。Beer則にも化学種・波長・散乱等の成立条件がある。

\paragraph{問題の要旨} 分析信号と水銀温度計・天秤の環境依存を比較する。

\Needspace{4\baselineskip}\paragraph{確認した条件・記述（独自要約）}\mbox{}\par

\noindent\textbf{原典確認：} S7\_049の26原典画像で全五肢・否定・添字を確認。\par

\noindent\textbf{対象：} 分析信号と水銀温度計・天秤の環境依存を比較する。\par

\paragraph{問われる核心} 分析信号と水銀温度計・天秤の環境依存を比較する。

\paragraph{必要知識} 天秤で下向きの有効力は概略mg−ρair Vg。空気密度は温度・圧力・湿度に依存し、試料と分銅の体積差によって浮力補正が必要となる。全く変わらない5を選ぶ。Beer則A=−log10 T=εclなのでH29のlog透過度も負の比例であり、R02の吸光度も正の比例となる。

\paragraph{考え方・正答を選ぶ理由（査読待ち初稿）} 天秤で下向きの有効力は概略mg−ρair Vg。空気密度は温度・圧力・湿度に依存し、試料と分銅の体積差によって浮力補正が必要となる。全く変わらない5を選ぶ。Beer則A=−log10 T=εclなのでH29のlog透過度も負の比例であり、R02の吸光度も正の比例となる。

\Needspace{4\baselineskip}\paragraph{選択肢の検討（査読待ち初稿）}\mbox{}\par

\noindent\textbf{1：通常の条件で適切} 炎中の原子励起と発光。高い状態から低い状態への遷移による特性発光を測る。\par

\noindent\textbf{2：通常の条件で適切} 吸光度の濃度比例。固定波長・光路長でBeer則A=−log10T=εcl。logTも係数の符号が負の比例関係となる。\par

\noindent\textbf{3：条件付き・解釈留保} 水銀柱全部が同温度なら正しい。全浸没用温度計ではその条件に対応する。一方、部分浸没型は指定浸没位置と露出柱の校正条件があり、全型に無条件の規則ではない。\par

\noindent\textbf{4：通常の条件で適切} XRF波長が化学結合状態で変わり得る。内殻の信号にも周辺電子状態に応じた小さいエネルギー差があり、高分解能では化学シフトを扱える。\par

\noindent\textbf{5：不適切} 天秤は温度・圧力・湿度で変化しない。空気密度と浮力が変わる。試料と分銅の密度差があると見かけの比較が変わる。\par

\paragraph{再利用する考え方} 測定対象と分離対象、定義の分母、絶対量と相対校正、全称と成立条件を分けて判断する。

\paragraph{共通理論} \hyperref[theory:analysis:metrology]{第\ref*{theory:analysis:metrology}節}

\par 参考：\cite{analysis49_beer_ref,analysis49_thermo_ref,analysis49_buoy_ref,analysis49_sxes_ref}。

\paragraph{関連問題} \hyperref[q:h29-reg-applied-12]{H29-REG-APPLIED-12}

\par 難易度 4（校準初稿）、優先度未評価。

\section{高分子分子量の相対GPC測定と浸透圧による数平均を比較する。}

\subsection{平成25年度 問16}
\label{q:h25-reg-applied-16}

\noindent\textbf{原典：} sources/H25/17\_applied\_science.pdf、PDF物理7頁。

\textbf{公式正答：} 3（single）。

\par\textbf{独立解答と公式正答に相違あり：} 独立解答 1|3。公式3を保持。初回独立1|3を保存。1も校正されたGPCによる分子量測定として成立し、単一正答の排他性は未解決。

\paragraph{問題の要旨} 高分子分子量の相対GPC測定と浸透圧による数平均を比較する。

\Needspace{4\baselineskip}\paragraph{確認した条件・記述（独自要約）}\mbox{}\par

\noindent\textbf{原典確認：} S7\_049の26原典画像で全五肢・否定・添字を確認。\par

\noindent\textbf{対象：} 高分子分子量の相対GPC測定と浸透圧による数平均を比較する。\par

\paragraph{問われる核心} 高分子分子量の相対GPC測定と浸透圧による数平均を比較する。

\paragraph{必要知識} 質量濃度c、数平均分子量Mnなら希薄溶液でΠ/c=\allowbreak{}RT/Mn+O(c)となる。濃度を変えてゼロへ外挿すれば切片からMnを得るので公式は3。一方、1のGPCも標準試料による相対校正で分子量を測定する実用法である。絶対測定の優位性で3を説明できるが、原典の一般的な問いから1を排除する十分な限定はない。

\paragraph{考え方・正答を選ぶ理由（査読待ち初稿）} 質量濃度c、数平均分子量Mnなら希薄溶液でΠ/c=\allowbreak{}RT/Mn+O(c)となる。濃度を変えてゼロへ外挿すれば切片からMnを得るので公式は3。一方、1のGPCも標準試料による相対校正で分子量を測定する実用法である。絶対測定の優位性で3を説明できるが、原典の一般的な問いから1を排除する十分な限定はない。

\Needspace{4\baselineskip}\paragraph{選択肢の検討（査読待ち初稿）}\mbox{}\par

\noindent\textbf{1：条件付き・解釈留保} ゲルろ過の溶離時間。サイズ排除クロマトグラフィーとして標準物質で校正すれば分子量と分布を測定できる。原典に絶対測定という限定はなく、一般の分子量測定法としても成立する。\par

\noindent\textbf{2：不適切} 中和滴定から解離度。示されたpH・解離定数の式は電離状態を調べる。末端基の数など追加情報なしに一般の高分子分子量を与えない。\par

\noindent\textbf{3：通常の条件で適切} 半透膜による浸透圧。希薄溶液の粒子数に依存する浸透圧から数平均分子量を求める。濃度ゼロへの外挿で相互作用の寄与を減らす。\par

\noindent\textbf{4：不適切} ガラス転移点の測定。分子量と相関することはあるが、化学構造・可塑剤等にも依存し、記述だけでは一意に分子量が決まらない。\par

\noindent\textbf{5：不適切} 円偏光二色性。左右円偏光に対する吸収差を測り、立体構造や二次構造を調べる。一般の分子量直接測定ではない。\par

\paragraph{再利用する考え方} 測定対象と分離対象、定義の分母、絶対量と相対校正、全称と成立条件を分けて判断する。

\paragraph{共通理論} \hyperref[theory:analysis:molar]{第\ref*{theory:analysis:molar}節}

\par 参考：\cite{analysis49_gpc_ref}。

\par 難易度 4（校準初稿）、優先度未評価。

\section{成分の分離とNMRによる観測を区別する。}

\subsection{平成25年度 問15}
\label{q:h25-reg-applied-15}

\noindent\textbf{原典：} sources/H25/17\_applied\_science.pdf、PDF物理6頁。

\textbf{公式正答：} 1（single）。

\par\textbf{検討未確定：} 質量分析でのイオン選別と、中性試料の調製・分取という分離の意味は同一でない。原典の広い成分分離の比較として公式のNMR選択を説明する。

\paragraph{問題の要旨} 成分の分離とNMRによる観測を区別する。

\Needspace{4\baselineskip}\paragraph{確認した条件・記述（独自要約）}\mbox{}\par

\noindent\textbf{原典確認：} S7\_049の26原典画像で全五肢・否定・添字を確認。\par

\noindent\textbf{対象：} 成分の分離とNMRによる観測を区別する。\par

\paragraph{問われる核心} 成分の分離とNMRによる観測を区別する。

\paragraph{必要知識} NMRは核の磁気共鳴を観測する分析法で、通常は成分を位置・経路別に分ける操作を行わない。遠心分離・電気泳動・HPLCは移動速度や保持の違いを利用する。MSは装置内でイオンをm/zで選別するという意味で分離を含む。候補番号の入替えを確認し、NMRを選ぶ。

\paragraph{考え方・正答を選ぶ理由（査読待ち初稿）} NMRは核の磁気共鳴を観測する分析法で、通常は成分を位置・経路別に分ける操作を行わない。遠心分離・電気泳動・HPLCは移動速度や保持の違いを利用する。MSは装置内でイオンをm/zで選別するという意味で分離を含む。候補番号の入替えを確認し、NMRを選ぶ。

\Needspace{4\baselineskip}\paragraph{選択肢の検討（査読待ち初稿）}\mbox{}\par

\noindent\textbf{1：不適切} 核磁気共鳴。共鳴周波数・化学シフト等を観測して構造を分析する。通常の装置は混合物を成分別に回収する分離器ではない。\par

\noindent\textbf{2：条件付き・解釈留保} 質量分析。イオンを質量電荷比で選別して測定する装置内の分離を含む。一方、元の中性混合物を大量に分取する手法とは異なる。\par

\noindent\textbf{3：通常の条件で適切} 遠心分離。回転場で密度・大きさに応じた沈降速度差を利用する。\par

\noindent\textbf{4：通常の条件で適切} 電気泳動。電荷と抵抗に依存する移動度差を利用して成分を分ける。\par

\noindent\textbf{5：通常の条件で適切} 高速液体クロマトグラフィー。固定相と移動相への分配等により保持時間が異なる成分を分ける。\par

\paragraph{再利用する考え方} 測定対象と分離対象、定義の分母、絶対量と相対校正、全称と成立条件を分けて判断する。

\paragraph{共通理論} \hyperref[theory:analysis:separation]{第\ref*{theory:analysis:separation}節}

\par 参考：\cite{analysis49_ms_ref}。

\paragraph{関連問題} \hyperref[q:h28-reg-applied-16]{H28-REG-APPLIED-16}

\par 難易度 4（校準初稿）、優先度未評価。

\subsection{平成28年度 問16}
\label{q:h28-reg-applied-16}

\noindent\textbf{原典：} sources/H28/17\_applied\_science.pdf、PDF物理6頁。

\textbf{公式正答：} 1（single）。

\par\textbf{検討未確定：} 質量分析でのイオン選別と、中性試料の調製・分取という分離の意味は同一でない。原典の広い成分分離の比較として公式のNMR選択を説明する。

\paragraph{問題の要旨} 成分の分離とNMRによる観測を区別する。

\Needspace{4\baselineskip}\paragraph{確認した条件・記述（独自要約）}\mbox{}\par

\noindent\textbf{原典確認：} S7\_049の26原典画像で全五肢・否定・添字を確認。\par

\noindent\textbf{対象：} 成分の分離とNMRによる観測を区別する。\par

\paragraph{問われる核心} 成分の分離とNMRによる観測を区別する。

\paragraph{必要知識} NMRは核の磁気共鳴を観測する分析法で、通常は成分を位置・経路別に分ける操作を行わない。遠心分離・電気泳動・HPLCは移動速度や保持の違いを利用する。MSは装置内でイオンをm/zで選別するという意味で分離を含む。候補番号の入替えを確認し、NMRを選ぶ。

\paragraph{考え方・正答を選ぶ理由（査読待ち初稿）} NMRは核の磁気共鳴を観測する分析法で、通常は成分を位置・経路別に分ける操作を行わない。遠心分離・電気泳動・HPLCは移動速度や保持の違いを利用する。MSは装置内でイオンをm/zで選別するという意味で分離を含む。候補番号の入替えを確認し、NMRを選ぶ。

\Needspace{4\baselineskip}\paragraph{選択肢の検討（査読待ち初稿）}\mbox{}\par

\noindent\textbf{1：不適切} 核磁気共鳴。共鳴周波数・化学シフト等を観測して構造を分析する。通常の装置は混合物を成分別に回収する分離器ではない。\par

\noindent\textbf{2：条件付き・解釈留保} 質量分析。イオンを質量電荷比で選別して測定する装置内の分離を含む。一方、元の中性混合物を大量に分取する手法とは異なる。\par

\noindent\textbf{3：通常の条件で適切} 遠心分離。回転場で密度・大きさに応じた沈降速度差を利用する。\par

\noindent\textbf{4：通常の条件で適切} 電気泳動。電荷と抵抗に依存する移動度差を利用して成分を分ける。\par

\noindent\textbf{5：通常の条件で適切} 高速液体クロマトグラフィー。固定相と移動相への分配等により保持時間が異なる成分を分ける。\par

\paragraph{再利用する考え方} 測定対象と分離対象、定義の分母、絶対量と相対校正、全称と成立条件を分けて判断する。

\paragraph{共通理論} \hyperref[theory:analysis:separation]{第\ref*{theory:analysis:separation}節}

\par 参考：\cite{analysis49_ms_ref}。

\paragraph{関連問題} \hyperref[q:h25-reg-applied-15]{H25-REG-APPLIED-15}, \hyperref[q:r01-reg-applied-12]{R01-REG-APPLIED-12}

\par 難易度 4（校準初稿）、優先度未評価。

\subsection{令和元年度 問12}
\label{q:r01-reg-applied-12}

\noindent\textbf{原典：} sources/R01/17\_applied\_science.pdf、PDF物理4頁。

\textbf{公式正答：} 4（single）。

\par\textbf{検討未確定：} 質量分析でのイオン選別と、中性試料の調製・分取という分離の意味は同一でない。原典の広い成分分離の比較として公式のNMR選択を説明する。

\paragraph{問題の要旨} 成分の分離とNMRによる観測を区別する。

\Needspace{4\baselineskip}\paragraph{確認した条件・記述（独自要約）}\mbox{}\par

\noindent\textbf{原典確認：} S7\_049の26原典画像で全五肢・否定・添字を確認。\par

\noindent\textbf{対象：} 成分の分離とNMRによる観測を区別する。\par

\paragraph{問われる核心} 成分の分離とNMRによる観測を区別する。

\paragraph{必要知識} NMRは核の磁気共鳴を観測する分析法で、通常は成分を位置・経路別に分ける操作を行わない。遠心分離・電気泳動・HPLCは移動速度や保持の違いを利用する。MSは装置内でイオンをm/zで選別するという意味で分離を含む。候補番号の入替えを確認し、NMRを選ぶ。

\paragraph{考え方・正答を選ぶ理由（査読待ち初稿）} NMRは核の磁気共鳴を観測する分析法で、通常は成分を位置・経路別に分ける操作を行わない。遠心分離・電気泳動・HPLCは移動速度や保持の違いを利用する。MSは装置内でイオンをm/zで選別するという意味で分離を含む。候補番号の入替えを確認し、NMRを選ぶ。

\Needspace{4\baselineskip}\paragraph{選択肢の検討（査読待ち初稿）}\mbox{}\par

\noindent\textbf{1：条件付き・解釈留保} 質量分析。イオンを質量電荷比で選別して測定する装置内の分離を含む。一方、元の中性混合物を大量に分取する手法とは異なる。\par

\noindent\textbf{2：通常の条件で適切} 電気泳動。電荷と抵抗に依存する移動度差を利用して成分を分ける。\par

\noindent\textbf{3：通常の条件で適切} 遠心分離。回転場で密度・大きさに応じた沈降速度差を利用する。\par

\noindent\textbf{4：不適切} 核磁気共鳴。共鳴周波数・化学シフト等を観測して構造を分析する。通常の装置は混合物を成分別に回収する分離器ではない。\par

\noindent\textbf{5：通常の条件で適切} 高速液体クロマトグラフィー。固定相と移動相への分配等により保持時間が異なる成分を分ける。\par

\paragraph{再利用する考え方} 測定対象と分離対象、定義の分母、絶対量と相対校正、全称と成立条件を分けて判断する。

\paragraph{共通理論} \hyperref[theory:analysis:separation]{第\ref*{theory:analysis:separation}節}

\par 参考：\cite{analysis49_ms_ref}。

\paragraph{関連問題} \hyperref[q:h28-reg-applied-16]{H28-REG-APPLIED-16}

\par 難易度 4（校準初稿）、優先度未評価。

\section{多元素同時分析に向く方法を比較する。}

\subsection{平成24年度 問16}
\label{q:h24-reg-applied-16}

\noindent\textbf{原典：} sources/H24/17\_applied\_science.pdf、PDF物理5頁。

\textbf{公式正答：} 5（single）。

\paragraph{問題の要旨} 多元素同時分析に向く方法を比較する。

\Needspace{4\baselineskip}\paragraph{確認した条件・記述（独自要約）}\mbox{}\par

\noindent\textbf{原典確認：} S7\_049の26原典画像で全五肢・否定・添字を確認。\par

\noindent\textbf{対象：} 多元素同時分析に向く方法を比較する。\par

\paragraph{問われる核心} 多元素同時分析に向く方法を比較する。

\paragraph{必要知識} 元素の種類を区別する観測量を探す。特性X線のエネルギー又は波長を検出するXRFなら、複数元素のピークを識別できる。各滴定の総当量だけでは混合元素を一度に分解できない。装置方式による同時・逐次検出の差はあるが、五方法の比較では5を選ぶ。

\paragraph{考え方・正答を選ぶ理由（査読待ち初稿）} 元素の種類を区別する観測量を探す。特性X線のエネルギー又は波長を検出するXRFなら、複数元素のピークを識別できる。各滴定の総当量だけでは混合元素を一度に分解できない。装置方式による同時・逐次検出の差はあるが、五方法の比較では5を選ぶ。

\Needspace{4\baselineskip}\paragraph{選択肢の検討（査読待ち初稿）}\mbox{}\par

\noindent\textbf{1：不適切} 酸化還元滴定。当量関係を使う濃度測定で、複数元素の特性線を一度に識別する方法ではない。\par

\noindent\textbf{2：不適切} 赤外吸収。主に振動と官能基を調べる。通常の多元素同時定量の第一選択ではない。\par

\noindent\textbf{3：不適切} キレート滴定。金属イオンと錯形成する試薬の当量を使い、各元素を同時に独立識別する方式ではない。\par

\noindent\textbf{4：不適切} 吸光光度。特定の吸収帯や呈色反応を用いる。ここでの全元素比較ではXRFが適する。\par

\noindent\textbf{5：通常の条件で適切} 蛍光X線。元素ごとの内殻遷移エネルギーを分け、複数の元素の信号を一分析で得られる。\par

\paragraph{再利用する考え方} 測定対象と分離対象、定義の分母、絶対量と相対校正、全称と成立条件を分けて判断する。

\paragraph{共通理論} \hyperref[theory:analysis:xray]{第\ref*{theory:analysis:xray}節}

\par 参考：\cite{analysis49_xrf_ref}。

\par 難易度 3（校準初稿）、優先度未評価。

\subsection{令和3年度 問16}
\label{q:r03-reg-applied-16}

\noindent\textbf{原典：} sources/R03/17\_applied\_science.pdf、PDF物理6頁。

\textbf{公式正答：} 3（single）。

\paragraph{問題の要旨} 内殻光電子の化学シフトと核・電子スピンを比較する。

\Needspace{4\baselineskip}\paragraph{確認した条件・記述（独自要約）}\mbox{}\par

\noindent\textbf{原典確認：} S7\_049の26原典画像で全五肢・否定・添字を確認。\par

\noindent\textbf{対象：} 内殻光電子の化学シフトと核・電子スピンを比較する。\par

\paragraph{問われる核心} 内殻光電子の化学シフトと核・電子スピンを比較する。

\paragraph{必要知識} 内殻の信号だから結合状態に無関係、という推論を避ける。XPSの結合エネルギーには化学シフトがあり3が不適切。励起X線はE=\allowbreak{}hc/λなので短波長ほど高エネルギーで、吸収端を超える内殻励起が蛍光X線を生む。

\paragraph{考え方・正答を選ぶ理由（査読待ち初稿）} 内殻の信号だから結合状態に無関係、という推論を避ける。XPSの結合エネルギーには化学シフトがあり3が不適切。励起X線はE=\allowbreak{}hc/λなので短波長ほど高エネルギーで、吸収端を超える内殻励起が蛍光X線を生む。

\Needspace{4\baselineskip}\paragraph{選択肢の検討（査読待ち初稿）}\mbox{}\par

\noindent\textbf{1：通常の条件で適切} ESRと不対電子。不対電子の磁気共鳴を測る。\par

\noindent\textbf{2：通常の条件で適切} 13Cはスピン1/2でNMR対象。13Cは核スピン1/2を持つ。豊富な12Cは核スピン0なので同じ測定対象ではない。\par

\noindent\textbf{3：不適切} XPSは結合状態に鈍感で化学分析不向き。周辺結合状態による内殻結合エネルギーの化学シフトを測定でき、元素と化学状態の分析に用いられる。\par

\noindent\textbf{4：通常の条件で適切} Ramanで振動・回転を調べる。非弾性散乱の周波数差が振動・回転遷移に対応する。\par

\noindent\textbf{5：通常の条件で適切} 吸収端より短いX線でXRF。光子エネルギーが内殻電離の閾値を超えると空孔が作られ、その後の遷移から特性X線を放射し得る。\par

\paragraph{再利用する考え方} 測定対象と分離対象、定義の分母、絶対量と相対校正、全称と成立条件を分けて判断する。

\paragraph{共通理論} \hyperref[theory:analysis:xray]{第\ref*{theory:analysis:xray}節}

\par 参考：\cite{analysis49_xps_ref,analysis49_esr_ref,atomic_ir_raman}。

\par 難易度 3（校準初稿）、優先度未評価。

\section{陽イオンの系統分析}

\subsection{平成26年度 問16}
\label{q:h26-reg-applied-16}

\noindent\textbf{原典：} sources/H26/17\_applied\_science.pdf、PDF物理7頁。

\textbf{公式正答：} 4（single）。

\par\textbf{検討未確定：} 酸性度・濃度・錯形成の条件は出題の操作列に依存し一般条件へ無限定に拡張しない。

\paragraph{問題の要旨} 塩化物、酸性下の硫化物、水酸化物を順に除いたろ液から、最後に硫化物として分離する陽イオンを特定する。

\Needspace{4\baselineskip}\paragraph{確認した条件・記述（独自要約）}\mbox{}\par

\noindent\textbf{I：} 塩酸を滴下し沈殿をろ別\par

\noindent\textbf{II：} Iのろ液中でチオアセトアミドCH3CSNH2の加水分解を行い、生成した沈殿をろ別\par

\noindent\textbf{III：} IIのろ液中へアンモニア水を滴下し、生成した沈殿をろ別\par

\noindent\textbf{IV：} IIIのろ液中へ硫化水素を滴下し、生成した沈殿をろ別\par

\paragraph{問われる核心} 分析化学

\paragraph{必要知識} 系統的定性分析、溶解度積、硫化水素の酸解離、pHと沈殿選択性

\paragraph{考え方・正答を選ぶ理由（査読待ち初稿）} 手順IでAg+をAgClとして除く。手順IIの酸性条件でCu2+が難溶性の硫化物となる。手順IIIでFe3+を水酸化物として除く。アンモニア性のろ液に硫化水素を加える手順IVではZn2+をZnSとして分離するので④。Na+はこれらの分離後も水溶液に残る。これは出題された古典的系統分析の条件に基づく説明で、実験条件一般への無条件な適用はしない。

\Needspace{4\baselineskip}\paragraph{選択肢の検討（査読待ち初稿）}\mbox{}\par

\noindent\textbf{1：不適切} Na+。操作I〜IVでの難溶沈殿として分離されず水溶液に残る。\par

\noindent\textbf{2：不適切} Ag+。Iの塩酸でAgClとして先に除く。\par

\noindent\textbf{3：不適切} Cu2+。IIの酸性硫化物の段階で先に除く。\par

\noindent\textbf{4：適切・選択対象} Zn2+。III後のアンモニア性ろ液にH2Sを加えるIVでZnSとして分離。\par

\noindent\textbf{5：不適切} Fe3+。IIIで水酸化物として先に除く。\par

\paragraph{再利用する考え方} 系統的定性分析、溶解度積、硫化水素の酸解離、pHと沈殿選択性

\paragraph{共通理論} \hyperref[theory:analytical:separation]{第\ref*{theory:analytical:separation}節}

\par 参考：\cite{qualitative_stanford}。

\par 難易度 3（校準初稿）、優先度未評価。

